% MAS1004 Assignment 1: long report.
%
% You fill this in with your agent, in your repository. AGENTS.md tells the
% agent how. Red text in the PDF is what still has to be replaced.
%
% - Pictures go in figures/. Copy them from results/, which is not in git.
% - Outputs go in outputs/. run.py saves what it prints in
%   results/<tag>_output.txt, and PLAN.md saves the other outputs in results/
%   with `| tee results/...`. Copy them here. \outputfile{outputs/check.txt}
%   prints a file exactly as it is. Every number in this report has to come
%   from such an output.
%
% To make the PDF: compress this folder into a zip, and on Overleaf choose
% New Project, Upload Project. Set Menu, Compiler to XeLaTeX, and Recompile.

\documentclass[11pt]{article}
\usepackage{mas1004}

\reporttitle{Assignment 1: Long report}
\studentname{Your name}
\studentid{Your student ID}

\begin{document}
\makeheader

% ---------------------------------------------------------------------------
\section*{Problem 1. Categories and images}

\subsection*{Why these categories}
I chose these four categories because they are types of quartz, characterised
by different colours. I'm interested in crystals, I own at least one crystal of
each category, and they are the ones I can usually distinguish easily.
They are also things I can photograph myself, which Problem 5 needs.

I checked the four against \texttt{imagenet\_classes.txt}, the 1,000 categories
ResNet18 already tells apart: none of the four is in it, the only entries that
share a word with them being "stone wall" and "ruddy turnstone", and nothing in
the list is a piece of polished quartz. So the last layer has to be replaced
and these four categories have to be learned from features that were trained on
something else.

\subsection*{What I expected the model to use}
Colour, above all, written down before training: purple is very different from
green, so I thought amethyst and aventurine would not get mixed up, and rose
quartz can be lighter in color, which makes it more similar to the light green
of aventurine, so I expected rose quartz and aventurine to be the pair that
does get mixed up. Shape I expected to be little help, because the pictures are
polished pieces, all of them roundish pebbles with no points to go by.

\subsection*{Images per category}
Asked \texttt{collect.py} for 150 images in each of four categories. Before any
cleaning, in \texttt{data/raw}:

\begin{center}\footnotesize
\begin{tabular}{lrr}
\toprule
class & asked for & came back \\
\midrule
amethyst & 150 & 134 \\
aventurine & 150 & 138 \\
citrine & 150 & 150 \\
rose\_quartz & 150 & 150 \\
\midrule
total & 600 & 572 \\
\bottomrule
\end{tabular}
\end{center}

Amethyst and aventurine came back short because \texttt{collect.py} skips a link
that does not open, is not an image, or is too small to use, and counts what it
skipped. The 134, 138, 150, 150 are the \texttt{downloaded} column of
\texttt{clean.py count} (\texttt{outputs/clean\_count.txt}), the counts as they
were before anything was removed; the 21 images for Problem 5, five in each class
and six rose quartz, are printed by \texttt{check.py} in the Appendix. Every
image is stored as \texttt{data/raw/<class name>/0001.jpg} onwards, and
\texttt{data/clean} was a copy of \texttt{data/raw} until \texttt{clean.py
remove} started moving files out of it.

% ---------------------------------------------------------------------------
\section*{Problem 2. Preparing an image}

\sidebyside{figures/prepare_original.jpg}{as downloaded: a rose quartz heart on a
            stand, with a sphere and two towers behind it}
           {figures/prepare_result.png}{what prepare\_image made of it: the middle
            224 \(\times\) 224 square, with the heart filling it}

\texttt{prepare\_image} resizes the shorter side to 256 and keeps the middle
224 \(\times\) 224 square. Measured with \texttt{src/show\_prepare.py} on
\texttt{data/new\_images/aventurine/IMG\_3068.JPG}, one of the photographs taken
by hand for this assignment, 4032 \(\times\) 3024:

\begin{center}\footnotesize
\begin{tabular}{ll}
\toprule
step & result \\
\midrule
on disk & 4032 \(\times\) 3024 \\
after Resize(256) & 341 \(\times\) 256 \\
square kept & x 58..282, y 16..240 \\
cut off, in the 341 \(\times\) 256 picture & 58 left, 59 right, 16 top, 16 bottom \\
cut off, back in the original photo & 686 left, 698 right, 189 top, 189 bottom \\
kept & 57.5\% of the area \\
\bottomrule
\end{tabular}
\end{center}

It throws away 34.3\% of the width and 12.5\% of the height: 686 pixels on each
side of that photo, even though the stone is nowhere near those edges. That the
rectangle drawn on the picture is the rectangle the model gets is measured, not
assumed: the same box through ToTensor and Normalize differs from
\texttt{prepare\_image} by 0.0e+00, and \texttt{prepare\_image} differs from
torchvision's own ImageNet recipe for ResNet18 by 0.0e+00.

An image where the crop cut off part of the thing I care about:
\texttt{data/new\_images/rose\_quartz/9dacd40a-f98e-4312-a62d-78c4fd3a0a4b.JPG},
900 \(\times\) 1600. Resize gives 256 \(\times\) 455 and the square keeps
y 116..340, so 116 rows go from above and 115 from below: 50.8\% of the height,
408 and 404 pixels of the original, and only 43.1\% of the area survives. The
stone reaches the top of the square, so its edge is thrown away.

A second case, found with \texttt{show\_prepare.py --scan data/new\_images
--sheet}: the contact sheet draws the 15 crops that throw most away, and
\texttt{data/new\_images/rose\_quartz/ea7f4b76-efcc-4973-814e-a7afe9273198.JPG},
1500 \(\times\) 2000, keeps 57.5\% of its area while the point of the heart sits
below the bottom edge of the square: the tip of the stone never reaches the
model. The sheet found four pieces the crop cuts through, all of them in these
21 photos.

\picturefile{figures/prepare_cropbox.png}
\picturefile{figures/prepare_side_by_side.png}
\picturefile[0.9\linewidth]{figures/prepare_sheet.png}

Four of my 21 photos are stored sideways: \texttt{amethyst/IMG\_3200},
\texttt{aventurine/IMG\_3196}, \texttt{rose\_quartz/IMG\_3069} and
\texttt{rose\_quartz/IMG\_3194}. The phone saved them lying on their side, with a
note in the file that says ``turn this photo before showing it''.
\texttt{prepare\_image} does not read that note, so in Python the model sees these
four photos turned on their side, while the web page does read it and shows them
upright. For these four photos, Python and the page therefore do not see quite
the same picture.
% ---------------------------------------------------------------------------
\section*{Problem 3. First training runs}

\begin{center}\footnotesize
\begin{tabular}{llllrrrr}
\toprule
tag & start & trained & epochs & lr & trainable parameters & train acc. & test acc. \\
\midrule
run & ImageNet & all layers & 10 & 0.0001 & 11,178,564 & 100.0\% & 92.2\% \\
scratch & random & all layers & 10 & 0.0001 & 11,178,564 & 99.6\% & 82.6\% \\
frozen & ImageNet & last layer & 10 & 0.001 & 2,052 & 97.2\% & 88.7\% \\
lr1e-3 & ImageNet & all layers & 10 & 0.001 & 11,178,564 & 96.7\% & 84.3\% \\
clean & ImageNet & all layers & 10 & 0.0001 & 11,178,564 & 100.0\% & 100.0\% \\
\bottomrule
\end{tabular}
\end{center}

The first four rows are comparable with each other: they are the same code on the
same \texttt{data/clean}, which at that point was still identical to
\texttt{data/raw}, so all four test accuracies are measured on the same 115
images, 457 train and 115 test out of 572. The fifth row, \texttt{clean}, is not
comparable with them: it is the plain \texttt{run} settings again (ImageNet
start, all layers, 10 epochs, lr 1e-4) on the same collection after the 198
removals of Problem 4, so 374 images and a different split, 299 train and 75
test. Only the data changed, which is what makes it a before-and-after pair with
\texttt{run}, and it is the pair Problems 4 and 5 compare.

One test image is 0.87\% of 115. \texttt{run.py} seeds the split but not the
shuffling or the new last layer's initial weights, so the same command does not
print the same number twice: running the first three commands again moved test
accuracy by +2.6 (+3 images), +1.7 (+2) and +3.5 (+4) points. Differences below
about four test images are therefore not evidence of anything, which makes
\texttt{run} (106/115) and \texttt{frozen} (102/115) a tie, and pretraining
against no pretraining (106 against 95) not one.

The first attempt, the same three commands run once before, with the same split
and only a warning in \texttt{train.py} that was fixed afterwards (the loss was
turned into a number without \texttt{.detach()}, which changes what is logged,
not what is trained), printed 89.6\%, 80.9\% and 85.2\% test accuracy
(103, 93 and 98 of 115). The table uses the second attempt. Its output, as
printed on Colab:

\outputfile{outputs/first_attempt_output.txt}

\picturefile{figures/run_curves.png}

The train loss in the \texttt{run} curve goes down from 0.8494 at epoch 1 to
0.0054 at epoch 10, with one small rise at epoch 4 (0.0485 to 0.0615). From
epoch 5 on it is below 0.02 and keeps going down only slowly. Test loss is a
different story. It falls to 0.2505 at epoch 4, rises to 0.3218 at epoch 9 and
ends at 0.2822, and test accuracy stops improving at epoch 4 as well: 92.2\% at
epoch 4, 90.4\% at epoch 5, and 92.2\% again for epochs 6 to 10. Epochs 5 to 10 bought nothing on
pictures the model had not seen. In the \texttt{scratch} curve test loss never
settles below about 0.52 and jumps to 1.2557 at epoch 7, while train accuracy keeps rising, to 98.9\% at
epoch 10: that is a network fitting its training set.

100.0\% train against 92.2\% test is a gap of 7.8 points, which is all 457
training images right against 9 of 115 test images wrong. It means the model can
memorise 457 pictures of four visually distinct stones, and that the test number
is the one that says something about pictures it has not seen.

\begin{center}\picturefile[0.5\linewidth]{figures/run_confusion.png}\end{center}

The nine wrong test images are all in the matrix below, printed from
\texttt{results/run\_model.pt} by \texttt{src/show\_confusion.py --tag run} (the
same split, seed 0, 115 pictures, so the picture \texttt{run\_confusion.png} and
these numbers agree):

\begin{center}\footnotesize
\begin{verbatim}
                 amethyst   aventurine      citrine  rose_quartz      total
amethyst                23            2            0            2         27
aventurine               0           25            1            2         28
citrine                  0            0           30            0         30
rose_quartz              2            0            0           28         30
                        25           27           31           32        115
\end{verbatim}
\end{center}

Recall per class: amethyst 23/27 (85\%), aventurine 25/28 (89\%), citrine 30/30
(100\%), rose\_quartz 28/30 (93\%). No single direction is ahead: 2 is the
largest count and four directions reach it, so which pair looks "most" mixed up
depends on whether the two directions of a pair are added. Added, the two biggest
pairs both have rose quartz in them, amethyst/rose\_quartz 4 and
aventurine/rose\_quartz 4, against amethyst/aventurine 2 and aventurine/citrine 1.
Citrine is never the true class of a mistake, and it ends with one more
prediction than it has pictures, 31 against 30.

My answer to "which two categories, and did it surprise you":

\begin{quote}
The categories that got mixed up and surprised me was amethyst turning out as
aventurine twice. That is because purple is very different from green, so I
thought they wouldn't get mixed. Aventurine was mixed up for rose quartz also
twice but that surprised me less compared to the amethyst turning out as
aventurine, because rose quartz can be lighter in color, which makes it more
similar to the light green of aventurine. I wasn't surprised about amethyst and
rose quartz to get mixed up due to their similar colors.
\end{quote}

Amethyst with rose quartz and aventurine with rose quartz are tied as the pairs
it mixes up most, 4 of the 9 mistakes each through both directions.
Amethyst with rose quartz is the pair I was not surprised about.

Same code, same data, same epochs and learning rate: 92.2\% starting from
ImageNet against 82.6\% starting from random numbers, 11 test images apart, and
the \texttt{scratch} test loss rises while its train accuracy approaches 100\%.
The pretrained features already carry most of what separates coloured quartz; the
random network has to learn all of it from 457 photographs.

Training only the last layer, 2,052 of 11,178,564 parameters
(2,052 / 11,178,564 = 0.018\%, 5,445 times fewer), reaches 88.7\% against 92.2\%
for training everything, and takes 0.1 minutes instead of 0.3. That is within
the four-image noise of the best run, with a smaller train/test gap (97.2 against
100.0) and a smoother curve.

A fourth setting, \texttt{lr1e-3}, is the default run with the learning rate ten
times larger: 84.3\% against 92.2\%, and an erratic test loss, 16.8645 at epoch
1, 1.8993 at epoch 4 and 1.4221 at epoch 9, where the run at 0.0001 starts at 0.6067 and keeps
test loss under 0.33 from epoch 3 on. A learning rate big enough to pull one fresh layer
into place is too big when all 11 million parameters can move.

The curves of the other three runs, \texttt{scratch}, \texttt{frozen} and
\texttt{lr1e-3}, top to bottom:

\picturefile[0.8\linewidth]{figures/scratch_curves.png}
\picturefile[0.8\linewidth]{figures/frozen_curves.png}
\picturefile[0.8\linewidth]{figures/lr1e-3_curves.png}

\subsection*{Output of the runs}
\outputfile{outputs/run_output.txt}
\outputfile{outputs/scratch_output.txt}
\outputfile{outputs/frozen_output.txt}
\outputfile{outputs/lr1e-3_output.txt}

% ---------------------------------------------------------------------------
\section*{Problem 4. Cleaning}

\subsection*{My rule}
Written down before any \texttt{clean.py remove} ran, while
\texttt{data/removed/} was still empty, and not tidied afterwards:

\begin{quote}
My rules:

I remove images that don't show a crystal at all

I remove images that show drawings

I remove images that show jewellery where the crystal is not predominant

I remove images that has text and that take much space in the photo.

I remove images that don't show clearly the crystal or not in the right
direction.

I remove images that has more than one type of crystal

I remove images that show the wrong crystal.
\end{quote}

Seven clauses. \texttt{clean.py remove} takes one \texttt{--reason} per clause and
\texttt{clean.py count} groups by that string, so the shorter string written into
the log is named with each clause below.

\subsection*{Three kinds of junk}
Mostly unrelated objects: clothes, car accessories, logos, drawings, and with
them characters and the wrong crystal. Three kinds, each with one of the files
that went to \texttt{data/removed/}:

\begin{itemize}
\item not a crystal at all (clothes, car accessories, logos, characters), removed
      by "The image doesn't show a crystal at all": \texttt{aventurine/0102.jpg}
\item a drawing rather than a photograph, removed by "The image is a drawing":
      \texttt{amethyst/0118.jpg}
\item the wrong crystal, removed by "The image shows the wrong crystal":
      \texttt{aventurine/0110.jpg}
\end{itemize}

\picturefile[0.3\linewidth]{figures/junk_1.jpg}\hfill
\picturefile[0.3\linewidth]{figures/junk_2.jpg}\hfill
\picturefile[0.3\linewidth]{figures/junk_3.jpg}

\subsection*{What I removed}
\outputfile{outputs/clean_count.txt}

The suspect pass found 27 pairs of images above 0.95 similarity. 16 of those pairs
were settled by my own rule, because both halves were junk under one of the seven
clauses and had already gone: that is the whole amethyst tail, where
\texttt{0123}, \texttt{0124}, \texttt{0126}, \texttt{0127} and \texttt{0128} are
near copies of each other, so the run removed under "The image doesn't show a
crystal at all" took the duplicates with it. The other nine groups, 19 pictures,
were still in \texttt{data/clean} and I decided them one by one on the copies
sheet, keeping one picture of each group and removing the rest under one new
reason, "The image is a near copy of another picture":

\begin{center}\footnotesize
\begin{tabular}{ll}
\toprule
kept & the near copies removed \\
\midrule
\texttt{citrine/0100.jpg} & \texttt{citrine/0018.jpg} \\
\texttt{citrine/0149.jpg} & \texttt{citrine/0032.jpg} \\
\texttt{citrine/0006.jpg} & \texttt{citrine/0069.jpg} \\
\texttt{citrine/0050.jpg} & \texttt{citrine/0099.jpg} \\
\texttt{citrine/0083.jpg} & \texttt{citrine/0041.jpg} \\
\texttt{rose\_quartz/0144.jpg} & \texttt{rose\_quartz/0150.jpg} \\
\texttt{amethyst/0033.jpg} & \texttt{amethyst/0055.jpg} \\
\texttt{amethyst/0022.jpg} & \texttt{amethyst/0031.jpg} \\
\texttt{aventurine/0134.jpg} & \texttt{aventurine/0123.jpg}, \texttt{aventurine/0132.jpg} \\
\bottomrule
\end{tabular}
\end{center}

Ten images removed, one kept of each group, decided on 13 October 2026: 27 pairs
above 0.95, 16 of them already gone as junk by my own rule, nine groups left, one
picture kept from each and the other ten removed.

The suspect pass ran on all 572 images before any removal, so these numbers are
about the data as downloaded. Each sheet draws the 32 images of a class that look
least like the rest of their class, and \texttt{suspects.csv} carries a row for
every image with its rank, rank 1 being the least typical. Removals by rank:

\begin{center}\footnotesize
\begin{tabular}{lrrrr}
\toprule
class & removed by my rule & of those, rank 1--20 & rank 21--32 & rank 33 and later \\
\midrule
amethyst & 38 & 11 & 6 & 21 \\
aventurine & 81 & 15 & 9 & 57 \\
citrine & 12 & 4 & 0 & 8 \\
rose\_quartz & 57 & 15 & 6 & 36 \\
\bottomrule
\end{tabular}
\end{center}

The first 20 of each class went by: amethyst no crystal 5, jewellery 3, drawing 1,
unclear 1, wrong crystal 1; aventurine no crystal 9, wrong crystal 2, jewellery 1,
more than one 1, unclear 1, text 1; citrine text 2, more than one 2; rose\_quartz
no crystal 5, jewellery 4, more than one 4, wrong crystal 1, text 1.

Did the suspects include junk I had missed when I looked myself? Yes, some of
them were images I thought were ok when I first went through them.

Not every suspect was removed. When I made my rules, I also went through all the
downloaded images and listed the ones to remove next to the rule they broke, so
an image was removed because it broke a rule, not because it was a suspect. For
example, \texttt{amethyst/0107.jpg} was on the amethyst suspect sheet, which
gave it 8\% amethyst, but it did not break any of my rules, so it is still in
\texttt{data/clean}.

The images that look least like the rest of their class are mostly junk by my rule
too, most strongly in aventurine and rose\_quartz, which are also the two classes
my rule removed most from. Citrine is where the two views disagree most, 4 of 20
of the first suspects were removed, and aventurine is where I went furthest on my
own: 57 of its 81 removals sit at rank 33 or later, in pictures the sheet never
showed me. "Removed but not on the suspect list at all" is always 0 and says
nothing, because every one of the 572 images has a rank; the number that answers
"did you find junk it did not list" is that rank 33 and later column, 21, 57, 8
and 36.

\subsection*{Before and after cleaning}
\outputfile{outputs/compare.txt}

The \texttt{clean} run, its curves and its confusion matrix on the 75 test
images:

\picturefile[0.8\linewidth]{figures/clean_curves.png}
\begin{center}\picturefile[0.5\linewidth]{figures/clean_confusion.png}\end{center}

\outputfile{outputs/clean_output.txt}

Nothing went down, but the pair of test numbers is not as big as it looks. 92.2\%
was 9 wrong of 115; 100\% is 0 wrong of 75, so after cleaning one test image is
worth 1.3\% instead of 0.87\%, the 198 removals took 28 images out of the test set
(aventurine's share of it fell from 28 to 11), and the cleaned test set is the part of the
collection that survived my rule, so the unusual pictures that are hardest to
classify are mostly gone from it. Some of the 7.8 points is a better
model and some of it is a kinder exam, and these two numbers cannot tell those
apart. Nor can 100\% of 75 be anything more than 75 pictures out of 75:
\texttt{run.py} printed "no mistakes to plot, which is suspicious. Check your
split." for this run, because with nothing wrong on the test set there is no
mistake picture to draw, so this run has no \texttt{clean\_worst.png}. The
training set also fell from 457 to 299 images, a third smaller, and the test
accuracy went up anyway.

The fair comparison is the row that did not change: the same 21 photographs from
a different source, 18/21 before and 18/21 after, which is the pair of numbers
Problem 5 is built on. Cleaning moved the number of right answers by nothing at
all, and the next section shows that it did not leave the pictures alone either.

% ---------------------------------------------------------------------------
\section*{Problem 5. Images from a new source}

Photographs of crystals taken for this assignment, by me and by family members
and friends I asked for them, saved as JPEG because the phones' HEIC files
cannot be read by the code: 21 of them, five in
each class and six rose\_quartz, in \texttt{data/new\_images/<class>/}. One of
them, \texttt{data/new\_images/aventurine/IMG\_3068.JPG}, is the photograph
measured in Problem 2.

They cannot be among the downloads: they were never online, because they were
taken for this purpose, while the four categories were downloaded from a search engine,
and a search cannot return a photograph that never existed on the web. The
overlap check tests that instead of assuming it: every new image is described and
compared with all 572 downloaded files, including the 198 that moved to
\texttt{data/removed/}, and the closest any of the 21 comes to any of them is
0.85, under the 0.95 the tool calls a near copy. The tool prints 570 rather than
572 because it hashes every file and counts each picture once, and two of the 572
downloaded files are byte-for-byte copies of another file in the set.

\outputfile{outputs/overlap.txt}

\begin{center}\footnotesize
\begin{tabular}{lrr}
\toprule
run & test accuracy & new images \\
\midrule
\texttt{run} & 92.2\% & 85.7\% (18/21) \\
\texttt{clean} & 100.0\% & 85.7\% (18/21) \\
\bottomrule
\end{tabular}
\end{center}

The two test accuracies are measured on different pictures, because the split is
drawn after cleaning, 92.2\% of 115 and then 100.0\% of 75 (Problem 4). The new
images are the same 21 photographs for both rows and their total is the same too,
18 of 21 either way.

Not the same eighteen pictures. Loading both saved models and asking each of them
about all 21 new images (the file named by the first eight characters of its
name):

\begin{center}\footnotesize
\begin{tabular}{llll}
\toprule
new image & really & \texttt{run} says & \texttt{clean} says \\
\midrule
\texttt{amethyst/44db3aaf} & amethyst & \textbf{rose\_quartz} & amethyst \\
\texttt{amethyst/a95f9f93} & amethyst & \textbf{rose\_quartz} & amethyst \\
\texttt{rose\_quartz/IMG\_3069} & rose\_quartz & \textbf{aventurine} & rose\_quartz \\
\texttt{citrine/2788e4fa} & citrine & citrine & \textbf{amethyst} \\
\texttt{citrine/58f05f6d} & citrine & citrine & \textbf{rose\_quartz} \\
\texttt{citrine/6dc1dd3c} & citrine & citrine & \textbf{rose\_quartz} \\
\bottomrule
\end{tabular}
\end{center}

Those six pictures are the only ones where the two models disagree, three each
way, and no picture is wrong in both: 15 right in both, three wrong in each, so
\texttt{run} is wrong on 3/21 and \texttt{clean} is wrong on 3/21. Cleaning traded
three right answers for three wrong ones, one class for another:

\begin{center}\footnotesize
\begin{tabular}{lrr}
\toprule
class & \texttt{run} on the new images & \texttt{clean} on the new images \\
\midrule
amethyst & 3/5 & 5/5 \\
aventurine & 5/5 & 5/5 \\
citrine & 5/5 & 2/5 \\
rose\_quartz & 5/6 & 6/6 \\
\bottomrule
\end{tabular}
\end{center}

It is the same boundary, moved to the other side. Every pair \texttt{run} got
wrong on the downloaded test set had rose quartz in it, amethyst/rose\_quartz 4
and aventurine/rose\_quartz 4, its two biggest pairs, with citrine perfect at
30/30. On the new images it calls two amethysts rose quartz; after cleaning those
three answers are right, and three citrines have gone over to rose quartz and
amethyst in their place. My guess at why, marked as a guess and not as a
measurement: my rule took 83 of 138 aventurine pictures and 58 of 150
rose\_quartz, but only 17 of 150 citrine, so citrine is a larger share of a much
smaller collection and the pictures that held the citrine boundary are the ones
that went. A run on the same split as before would settle it; 21 pictures, one of
them worth 4.8\%, cannot.

\texttt{check.py} prints the five mistakes it was most sure about, and the
published model made three, so three is the whole list and not a selection. All
three are one small stone, in flat light, on a plain pale background:

\picturefile[0.32\linewidth]{figures/mistake_1_citrine_called_rose_quartz.jpg}\hfill
\picturefile[0.32\linewidth]{figures/mistake_2_citrine_called_rose_quartz.jpg}\hfill
\picturefile[0.32\linewidth]{figures/mistake_3_citrine_called_amethyst.jpg}

{\small left to right: said rose quartz 54\%, said rose quartz 51\%, said amethyst
32\%; all three really citrine.}

\begin{itemize}
\item \texttt{citrine/6dc1dd3c-0216-...JPG}, said rose\_quartz (54\%). One small
stone on grey speckled granite, and it is the pale end of citrine: milky white
with a faint yellow patch at one end, with the wet-looking grey paving around it
nearly as pale as the stone, so the whole frame is cream and grey. Nothing in the
picture is pink, but "pale stone, little colour" is what rose\_quartz looks like
in the downloads, and the flat light has taken the yellow away. Colour is the only
thing separating citrine from rose\_quartz in a picture like this, and the colour
is washed out.
\item \texttt{citrine/58f05f6d-77f2-...JPG}, said rose\_quartz (51\%). One amber
stone on white paper, underexposed, so the paper reads light grey and the amber
reads pale. 51\% is a coin toss, for the reason above.
\item \texttt{citrine/2788e4fa-a7ad-...JPG}, said amethyst (32\%). One pale yellow
stone on light grey paper, small in the frame, low contrast, most of the picture
empty background. 32\% out of four classes is close to guessing and there is
nothing purple in the picture: this one reads as an image the model had almost no
information from, rather than an image that fooled it in a particular way.
\end{itemize}

The three the earlier model got wrong instead are the three other wrong answers in
the table above: two amethysts called rose quartz, one a crystal deep violet over
half of it and white over the other half on pale pinkish paper, one a milky pale
lilac pendant with a silver chain on grey speckled granite; and one rose quartz
called aventurine, a pale lavender cabochon with a bright white highlight on a
grey surface, where aventurine in the downloads is also a smooth translucent
pebble, so the shape is no help. All six mistakes of the two models together are
the same shot: one small stone, flat light, a plain white or grey background, the
stone filling little of the frame, its colour washed out. The one citrine picture
the cleaned model does get right, \texttt{citrine/f16efa23}, has three stones,
warm light, saturated honey colour and hard shadows.

Compared with the downloads, my new images are less saturated, and their
colours are sometimes less vivid. The crystals are smaller, usually in the
middle of the photo. The light is darker in some of them. Some photos also show
my hand as the surface the crystal is lying on.

My guess at why cleaning did not make the model better on my new images: some
of the photos I removed because they broke my rules may have looked more like
my own photos, and the ones that remained look even more different from my
photos. Not all of the removed images were obvious junk. The images I removed are the
ones that broke my rules, plus the near copies I chose to remove, and some of
them were images I had thought were ok when I first went through them. Not every
suspect was removed: some of the images the tool listed as least like the rest
of their category did not break my rules, so I kept them. However, when I check the removed images, they do not necessarily look like my
photos. What I can see is that in the clean images the colours are more vivid
and the crystals are bigger in size than in my photos.

The suspect sheets (\texttt{outputs/suspects.csv}) show that citrine is
different from the other three categories. Each sheet shows the 32 downloaded
images of a category that look least like the rest of it. On the citrine sheet
the lowest probability of citrine is 53.1\%, 21 of the 32 are at 90\% or more,
and not one of them looks like another category. On the other sheets the lowest
probability is 0.03\% (amethyst), 0.55\% (aventurine) and 0.00\% (rose quartz),
and 13, 16 and 21 of their 32 are under 50\%. So the downloaded citrines all
look very much alike.

\picturefile{figures/suspects_citrine.png}

{\small The citrine suspect sheet, as \texttt{clean.py suspects} drew it.}

Even so, all three mistakes of the cleaned model on my
new images are citrines, while it got all 27 citrines in the downloaded test
set right.

My guess: the citrines in my photos are small, less vivid in colour and
sometimes paler, rather than having the full yellow or orange colour of the
citrines in the clean images, and in two of the three photos it got wrong the
light is darker.

% ---------------------------------------------------------------------------
\section*{Problem 6. The demo}

Address: \url{https://giuliag9.github.io/MAS1004-assignment1/}

I showed the page photos of four things that are none of my categories, and
it named one of my four stones every time, because it has no answer for ``none
of these'': an adapter was rose\_quartz (47.1\%, then amethyst 43.7\%), a
yellow highlighter was aventurine (65.3\%, then citrine only 24.9\%), a pink
mask product was rose\_quartz (92.2\%), and a beige mug was rose\_quartz
(79.4\%).

% ---------------------------------------------------------------------------
\section*{Problem 7. What the model looks at}

% Copy this subsection once for every guess.
\subsection*{Guess 1: colour separates amethyst from aventurine}

My guess, in my words as I gave it: ``I'll start from my original guess on what
the model would look at to distinguish the crystal. I mentioned I expect it to
look at colors. From my images, amethyst and aventurine were not mistaken in the
current model. So my first guess is that if we lower the saturation on amethyst
and aventurine the model will have more problems to distinguish them.''

The change: the saturation of all ten of my new amethyst and aventurine photos
lowered to 0\%, which leaves the light and dark, the shape and the texture of
each photo as they were and takes the colour away completely, so that I can
test the model with the absence of colour. Ten photos, five of each class, and
all ten originals are answered right by the page.

My requests to the agent, word for word:
\begin{quote}
make a copy of each image in data/new\_images/amethyst and
data/new\_images/aventurine with lowered saturation and save the copies,
respectively, in data/changed/less\_saturation\_amethyst and
data/changed/less\_saturation\_aventurine
\end{quote}
and, asked by how much:
\begin{quote}
to 0\% , so we can test with the absence of color.
\end{quote}

The agent wrote the change as a cell that I ran on Colab, where the images are:
Pillow's \texttt{ImageEnhance.Color(image).enhance(0.0)} on each photo,
written to \texttt{data/changed/less\_saturation\_<class>/} under the same file
name, with the originals left as they were. I uploaded every original and every
changed photo to my published page and read its first two answers off it:

\begin{center}\footnotesize
\begin{tabular}{lll}
\toprule
photo & original & saturation 0\% \\
\midrule
\texttt{amethyst/44db3aaf} & amethyst 96.7\%, rose quartz 2.4\% & amethyst 84.6\%, aventurine 7.9\% \\
\texttt{amethyst/a95f9f93} & amethyst 86.9\%, rose quartz 8.1\% & amethyst 77.0\%, aventurine 14.3\% \\
\texttt{amethyst/IMG_3072} & amethyst 80.2\%, rose quartz 16.7\% & \textbf{aventurine 45.5\%}, amethyst 36.4\% \\
\texttt{amethyst/IMG_3089} & amethyst 90.5\%, aventurine 5.2\% & amethyst 40.8\%, aventurine 30.6\% \\
\texttt{amethyst/IMG_3200} & amethyst 94.8\%, rose quartz 2.3\% & \textbf{aventurine 45.6\%}, amethyst 42.8\% \\
\texttt{aventurine/59ad40f0} & aventurine 99.3\%, amethyst 0.4\% & aventurine 62.9\%, amethyst 23.7\% \\
\texttt{aventurine/aba4d635} & aventurine 97.3\%, amethyst 1.7\% & aventurine 82.1\%, amethyst 12.5\% \\
\texttt{aventurine/IMG_3068} & aventurine 75.2\%, rose quartz 15.6\% & aventurine 51.5\%, rose quartz 33.7\% \\
\texttt{aventurine/IMG_3196} & aventurine 80.9\%, rose quartz 13.3\% & \textbf{rose quartz 47.1\%}, aventurine 44.6\% \\
\texttt{aventurine/IMG_3197} & aventurine 88.7\%, rose quartz 9.9\% & \textbf{rose quartz 62.3\%}, aventurine 34.6\% \\
\bottomrule
\end{tabular}
\end{center}

The originals are 10 of 10 right; with the colour taken away, 6 of 10 are still
right, and the probability of the right answer is lower for all ten. The four
that changed: two amethysts became aventurine, and two aventurines became rose
quartz.

% One \sidebyside for every image you changed. The pictures are scaled down for
% the report.
\sidebyside{figures/guess1_original_amethyst_44db3aaf.jpg}{amethyst/44db3aaf, original: amethyst 96.7\%, rose quartz 2.4\%}
           {figures/guess1_changed_amethyst_44db3aaf.jpg}{saturation 0\%: amethyst 84.6\%, aventurine 7.9\%}
\sidebyside{figures/guess1_original_amethyst_a95f9f93.jpg}{amethyst/a95f9f93, original: amethyst 86.9\%, rose quartz 8.1\%}
           {figures/guess1_changed_amethyst_a95f9f93.jpg}{saturation 0\%: amethyst 77.0\%, aventurine 14.3\%}
\sidebyside{figures/guess1_original_amethyst_IMG_3072.jpg}{amethyst/IMG_3072, original: amethyst 80.2\%, rose quartz 16.7\%}
           {figures/guess1_changed_amethyst_IMG_3072.jpg}{saturation 0\%: \textbf{aventurine 45.5\%}, amethyst 36.4\%}
\sidebyside{figures/guess1_original_amethyst_IMG_3089.jpg}{amethyst/IMG_3089, original: amethyst 90.5\%, aventurine 5.2\%}
           {figures/guess1_changed_amethyst_IMG_3089.jpg}{saturation 0\%: amethyst 40.8\%, aventurine 30.6\%}
\sidebyside{figures/guess1_original_amethyst_IMG_3200.jpg}{amethyst/IMG_3200, original: amethyst 94.8\%, rose quartz 2.3\%}
           {figures/guess1_changed_amethyst_IMG_3200.jpg}{saturation 0\%: \textbf{aventurine 45.6\%}, amethyst 42.8\%}
\sidebyside{figures/guess1_original_aventurine_59ad40f0.jpg}{aventurine/59ad40f0, original: aventurine 99.3\%, amethyst 0.4\%}
           {figures/guess1_changed_aventurine_59ad40f0.jpg}{saturation 0\%: aventurine 62.9\%, amethyst 23.7\%}
\sidebyside{figures/guess1_original_aventurine_aba4d635.jpg}{aventurine/aba4d635, original: aventurine 97.3\%, amethyst 1.7\%}
           {figures/guess1_changed_aventurine_aba4d635.jpg}{saturation 0\%: aventurine 82.1\%, amethyst 12.5\%}
\sidebyside{figures/guess1_original_aventurine_IMG_3068.jpg}{aventurine/IMG_3068, original: aventurine 75.2\%, rose quartz 15.6\%}
           {figures/guess1_changed_aventurine_IMG_3068.jpg}{saturation 0\%: aventurine 51.5\%, rose quartz 33.7\%}
\sidebyside{figures/guess1_original_aventurine_IMG_3196.jpg}{aventurine/IMG_3196, original: aventurine 80.9\%, rose quartz 13.3\%}
           {figures/guess1_changed_aventurine_IMG_3196.jpg}{saturation 0\%: \textbf{rose quartz 47.1\%}, aventurine 44.6\%}
\sidebyside{figures/guess1_original_aventurine_IMG_3197.jpg}{aventurine/IMG_3197, original: aventurine 88.7\%, rose quartz 9.9\%}
           {figures/guess1_changed_aventurine_IMG_3197.jpg}{saturation 0\%: \textbf{rose quartz 62.3\%}, aventurine 34.6\%}

What the results say about the guess, in my words: From these results I can
observe two things. First, the cases in which the 0\% saturation changed the
model's answer as I thought were only amethyst to aventurine, and not
aventurine to amethyst: the only two times the answer changed for aventurine,
it went to rose quartz and not to amethyst. Second, by looking at the changed
images I realised that the aventurine images that were called rose quartz at
0\% saturation were paler than the others, meaning that in the original the
colour was lighter, paler, less full. At the same time the amethyst images that
were called aventurine at 0\% saturation also showed some paler parts and
appeared less full, but still darker than the aventurine images that were
called rose quartz. Overall the answers change a bit after removing the colour,
but not as much as I expected, which makes me think that colour is considered,
but it is not the determining element for the model to tell the crystals
apart.

\subsection*{Guess 2: transparency, aventurine pale but dark}

My guess, in my words as I gave it: ``Given what I observed from the first
guess, color is not the determinant element for the model. I start to think
that transparency also is important when guessing the crystals. As I mentioned
in the first guess, 0\% saturation aventurine that were guessed to be rose quartz
were paler and less full in color, while the amethyst 0\% saturation that were
guessed to be aventurine also were less full, paler but overall darker. Which
makes me think that for the model aventurine is dark, but also pale, while rose
quartz is pale but lighter. If I think of the mistakes made, the first one was
of data/new\_images/citrine/6dc1dd3c-0216-45b8-abb7-214d633d280e.JPG being
guessed as rose quartz, and in fact the crystal in the photo is less full in
color, and lighter. To check this, the transparency of the images should be
lowered. In particular my guess to check if transparency plays a role in how
the model distinguishes the crystals, is that if we lower the transparency of
the images of amethyst we turned to 0\% saturation, since they're dark, they
would be guessed as aventurine (pale but dark).''

The change: a photo has no setting for how see-through the stone is, and the
transparency setting a picture file can have is removed by
\texttt{prepare\_image} when it converts the picture to RGB. So I chose to
convey less transparency by lowering the contrast to 50\%, on the whole picture,
on the three amethyst photos that were still called amethyst at 0\% saturation.
Lowering the contrast moves every part of the picture halfway towards its
average grey, so the crystal becomes paler and less full; I chose 50\% as a
compromise, to see the crystal paler without its dark parts getting too much
lighter.

My requests to the agent, word for word:
\begin{quote}
So make a copy of
data/changed/less\_saturation\_amethyst/44db3aaf-5cdf-4004-9255-868fefec64c3.JPG,
data/changed/less\_saturation\_amethyst/a95f9f93-5117-4b70-9e92-4c513bb5decc.JPG
and data/changed/less\_saturation\_amethyst/IMG\_3089.JPG with lower transparency
and save the copies in data/changed/less\_saturation\_less\_transparency\_amethyst
\end{quote}
and, after the agent explained that a picture has no transparency setting the
model would see:
\begin{quote}
Since there is no transparency setting to change in the images, I think the
best way to make the crystal paler and less full would be to lower contrast to
50\%, this way it's enough to make the image paler but keeps the darkess of the
crystal which in my guess is what determines aventurine (pale but dark). So make
a copy of
data/changed/less\_saturation\_amethyst/44db3aaf-5cdf-4004-9255-868fefec64c3.JPG,
data/changed/less\_saturation\_amethyst/a95f9f93-5117-4b70-9e92-4c513bb5decc.JPG
and data/changed/less\_saturation\_amethyst/IMG\_3089.JPG with contrast to 50\%
and save the copies in data/changed/less\_saturation\_less\_transparency\_amethyst.
Also maintain the whole picture of the images for this guess.
\end{quote}

The agent wrote the change as a cell that I ran on Colab: Pillow's
\texttt{ImageEnhance.Contrast(image).enhance(0.5)} on each of the three 0\%
saturation photos. The page's first two answers, next to the ones the same
photos got at 0\% saturation in Guess 1:

\begin{center}\footnotesize
\begin{tabular}{lll}
\toprule
photo & saturation 0\% & saturation 0\% and contrast 50\% \\
\midrule
\texttt{amethyst/44db3aaf} & amethyst 84.6\%, aventurine 7.9\% & \textbf{rose quartz 50.0\%}, amethyst 27.7\% \\
\texttt{amethyst/a95f9f93} & amethyst 77.0\%, aventurine 14.3\% & \textbf{rose quartz 46.2\%}, aventurine 27.8\% \\
\texttt{amethyst/IMG_3089} & amethyst 40.8\%, aventurine 30.6\% & \textbf{rose quartz 64.4\%}, aventurine 25.2\% \\
\bottomrule
\end{tabular}
\end{center}

All three changed their answer, all three to rose quartz and none to aventurine;
aventurine is second for two of them.

\sidebyside{figures/guess1_changed_amethyst_44db3aaf.jpg}{44db3aaf, saturation 0\%: amethyst 84.6\%, aventurine 7.9\%}
           {figures/guess2_changed_amethyst_44db3aaf.jpg}{and contrast 50\%: \textbf{rose quartz 50.0\%}, amethyst 27.7\%}
\sidebyside{figures/guess1_changed_amethyst_a95f9f93.jpg}{a95f9f93, saturation 0\%: amethyst 77.0\%, aventurine 14.3\%}
           {figures/guess2_changed_amethyst_a95f9f93.jpg}{and contrast 50\%: \textbf{rose quartz 46.2\%}, aventurine 27.8\%}
\sidebyside{figures/guess1_changed_amethyst_IMG_3089.jpg}{IMG_3089, saturation 0\%: amethyst 40.8\%, aventurine 30.6\%}
           {figures/guess2_changed_amethyst_IMG_3089.jpg}{and contrast 50\%: \textbf{rose quartz 64.4\%}, aventurine 25.2\%}

What the results say about the guess, in my words: These results make me think
that while lowering the contrast, the crystals also got lighter, and therefore
closer to rose quartz rather than aventurine. Lowering the contrast and keeping
the original darkness is difficult. However, these results seem to fit my idea
that rose quartz is pale and light. Even if it did not separate amethyst from
aventurine, transparency, which in this case was conveyed by lowering the
contrast, seems to be considered by the model: all three images got a different
answer from the one they had at 0\% saturation. It is also worth noticing that
the lower contrast changed all three answers, but it was applied to the 0\%
saturation images, not to the original ones. Overall, the results of the second
guess make me think of two things. First, if rose quartz is light and pale, and
the mistaken citrine is also light and pale, what does the model look at in the
image to tell them apart, and why was the citrine called rose quartz? Perhaps
transparency is not the answer. Second, there might be another factor that the
model considers to tell the crystals apart, or a combination of factors.

\subsection*{Guess 3: more colour would help the pale citrines}

My guess, in my words as I gave it: ``Given the previous guesses results, I
think color and transparency are in a way considered by the model, even if not
in the way I expected them. At this point is good to consider the case of the 3
mistakes of citrine of my images. Starting from the one I mentioned, citrine
6dc1dd3c-0216-45b8-abb7-214d633d280e.JPG was guessed as rose quartz, and these
crystals are generally light and pale. As I previously said, the images of
citrine in clean, are more vivid in color than my own photos, which I think is
an important element that makes the images in clean different from mine. In the
first guess I focused on color, and observed if its absence (so a 0\%
saturation) could impact the model's guessing, and overall, there were some
changes, even for the images that were guessed correctly (guessed with less
percentage). This makes me still want to keep color in consideration for a
guess. Therefore, it would be interesting to test the opposite of what I did in
the first guess, meaning instead of reducing color, enhancing it. Since both
citrine and rose quartz are pale, the aspect that is different in them is color
and them being mistaken could be linked to the color not being vivid. So, my
third guess is that if citrine has higher saturation the model would have less
problem recognizing it. If this guess turned to be true, it would confirm the
reason for the mistake in my images.''

What I expected, written before I saw the results: ``If my guess is right the
3 wrong will turn right and the other 2 right ones will stay right with higher
percentages.''

The change: 70\% more colour (saturation 170\% on the scale of Guess 1, where
100\% is the photo as it is), on the whole picture, on all five of my new
citrine photos: the three the model gets wrong and the two it gets right.

My requests to the agent, word for word:
\begin{quote}
So make a copy of each image in data/new\_images/citrine with higher saturation
(70\%) and save the copies in data/changed/saturated\_citrine
\end{quote}
and, asked whether 70\% meant 70\% more colour:
\begin{quote}
yes 70\% more color.
\end{quote}

The agent wrote the change as a cell that I ran on Colab: Pillow's
\texttt{ImageEnhance.Color(image).enhance(1.7)} on each photo. Raising the
saturation strengthens the colour that is there and cannot add colour to parts
that are white or grey. The page's first two answers:

\begin{center}\footnotesize
\begin{tabular}{lll}
\toprule
photo & original & saturation +70\% \\
\midrule
\texttt{citrine/6dc1dd3c} & rose quartz 52.8\%, aventurine 26.5\% & rose quartz 37.7\%, aventurine 31.5\% \\
\texttt{citrine/58f05f6d} & rose quartz 49.1\%, aventurine 22.4\% & rose quartz 54.7\%, aventurine 18.0\% \\
\texttt{citrine/2788e4fa} & rose quartz 31.0\%, amethyst 29.8\% & rose quartz 33.0\%, amethyst 25.1\% \\
\texttt{citrine/08f618d6} & citrine 89.5\%, aventurine 4.7\% & citrine 96.2\%, aventurine 2.1\% \\
\texttt{citrine/f16efa23} & citrine 59.9\%, amethyst 25.1\% & citrine 94.1\%, amethyst 4.1\% \\
\bottomrule
\end{tabular}
\end{center}

The two that were right stayed right and became more sure, \texttt{f16efa23}
from 59.9\% to 94.1\%. The three that were wrong stayed wrong, all three still
rose quartz, and citrine is not among the first two answers for any of them.

The page's numbers for the three originals are close to, but not the same as,
the ones \texttt{check.py} printed in Problem 5 (54\%, 51\% and 32\%), and for
\texttt{2788e4fa} the order of the first two even changes: \texttt{check.py}
said amethyst, the page says rose quartz 31.0\% with amethyst 29.8\% second. The
page shrinks a photo in the browser (\texttt{resizeShorterSide} in
\texttt{docs/app.js}) and Python shrinks it with Pillow, so the 224
\(\times\) 224 picture the model gets differs very slightly. The green badge
cannot see this, because its test pictures in \texttt{selftest.json} are already
224 \(\times\) 224. A photo the model is unsure about can change its answer with
a change too small to see.

\sidebyside{figures/guess3_original_citrine_6dc1dd3c.jpg}{citrine/6dc1dd3c, original: rose quartz 52.8\%, aventurine 26.5\%}
           {figures/guess3_changed_citrine_6dc1dd3c.jpg}{saturation +70\%: rose quartz 37.7\%, aventurine 31.5\%}
\sidebyside{figures/guess3_original_citrine_58f05f6d.jpg}{citrine/58f05f6d, original: rose quartz 49.1\%, aventurine 22.4\%}
           {figures/guess3_changed_citrine_58f05f6d.jpg}{saturation +70\%: rose quartz 54.7\%, aventurine 18.0\%}
\sidebyside{figures/guess3_original_citrine_2788e4fa.jpg}{citrine/2788e4fa, original: rose quartz 31.0\%, amethyst 29.8\%}
           {figures/guess3_changed_citrine_2788e4fa.jpg}{saturation +70\%: rose quartz 33.0\%, amethyst 25.1\%}
\sidebyside{figures/guess3_original_citrine_08f618d6.jpg}{citrine/08f618d6, original: citrine 89.5\%, aventurine 4.7\%}
           {figures/guess3_changed_citrine_08f618d6.jpg}{saturation +70\%: citrine 96.2\%, aventurine 2.1\%}
\sidebyside{figures/guess3_original_citrine_f16efa23.jpg}{citrine/f16efa23, original: citrine 59.9\%, amethyst 25.1\%}
           {figures/guess3_changed_citrine_f16efa23.jpg}{saturation +70\%: citrine 94.1\%, amethyst 4.1\%}

What the results say about the guess, in my words: Once again, the results do
not confirm my guess completely; in fact they only confirmed my guess for the
images of citrine that were already guessed right. However, the results give me
more insight about where I should put my attention in understanding what the
model looks at to distinguish the crystals, in particular what in those 3 images
of citrine is different from the others. As the saturation was increased, other
colours in the image also became more evident. In particular, it is possible to
notice that the background's colour is also more vivid. In the saturated image
6dc1dd3c-0216-45b8-abb7-214d633d280e.JPG the granite background shows more hues
that can resemble pink, but also a slight green. In the saturated image
58f05f6d-77f2-467e-b497-c4265bb43b8e.JPG, due to the type of light, the
background looks more pink too. The same happens in the saturated image
2788e4fa-a7ad-453a-b3ad-3984b54d6e2c.JPG. In fact, the results show that for the
saturated images 58f05f6d-77f2-467e-b497-c4265bb43b8e.JPG and
2788e4fa-a7ad-453a-b3ad-3984b54d6e2c.JPG the model is surer about its answer
``rose quartz'', while for 6dc1dd3c-0216-45b8-abb7-214d633d280e.JPG the first
answer, rose quartz, decreases to 37.7\% as aventurine rises to 31.5\%.
Therefore, the conclusion I draw from the results of the third guess is that the
colour of the background, also influenced by light, can influence the model's
answer in some situations.

\subsection*{Guess 4: the stone's size and the background}

My guess, in my words as I gave it: ``From the third guess I concluded that the
color of background can be a factor that influences the model. Given all the
previous guesses, I start to consider more and more the idea that the model uses
a combination of elements to guess the crystals. The images that were still
guessed wrong also had another thing in common: the crystals are small.
Additionally, the two citrine images that are guessed right are the only two
with more than one crystal in them. To check both background and size at the
same time, the images could be cut so that the crystal occupies the most part of
the image, and has little background. My guess would be that if we cut the
background from the images of citrine, making the crystal the main object in the
image, the model could guess them correctly.''

The change: the three citrines that were still wrong in Guess 3, in their
saturated versions, cut to a square around the stone, centred on it, its side
1.5 times the stone's longer side, so that the stone fills most of the picture
with a thin border of background. The square is the shape the model looks at,
so the page keeps all of the cut picture.

My request to the agent, word for word:
\begin{quote}
So make a copy of images
data/changed/saturated\_citrine/6dc1dd3c-0216-45b8-abb7-214d633d280e.JPG,
data/changed/saturated\_citrine/2788e4fa-a7ad-453a-b3ad-3984b54d6e2c.JPG and
data/changed/saturated\_citrine/58f05f6d-77f2-467e-b497-c4265bb43b8e.JPG cut the
background around the crystals and save the copies in
data/changed/saturated\_cut\_citrine. In this case still work on the saturated
version because I still think that saturation of the crystal plays a role.
\end{quote}

The agent found the stone in each picture, chose the square, showed me the three
cut pictures before they were made, and wrote the cut as a cell that I ran on
Colab. The page's first two answers, next to the ones the same saturated
pictures got in Guess 3:

\begin{center}\footnotesize
\begin{tabular}{lll}
\toprule
photo & saturated & saturated and cut \\
\midrule
\texttt{citrine/6dc1dd3c} & rose quartz 37.7\%, aventurine 31.5\% & \textbf{citrine 75.3\%}, aventurine 12.0\% \\
\texttt{citrine/58f05f6d} & rose quartz 54.7\%, aventurine 18.0\% & \textbf{citrine 99.6\%}, amethyst 0.3\% \\
\texttt{citrine/2788e4fa} & rose quartz 33.0\%, amethyst 25.1\% & \textbf{citrine 97.9\%}, amethyst 0.9\% \\
\bottomrule
\end{tabular}
\end{center}

All three changed from rose quartz to citrine, two of them above 97\%. The
pictures were both saturated and cut, so these three results alone cannot
separate the two changes; but in Guess 3 the same three pictures, saturated and
not cut, were all still rose quartz.

\sidebyside{figures/guess3_changed_citrine_6dc1dd3c.jpg}{6dc1dd3c, saturated: rose quartz 37.7\%, aventurine 31.5\%}
           {figures/guess4_changed_citrine_6dc1dd3c.jpg}{saturated and cut: citrine 75.3\%, aventurine 12.0\%}
\sidebyside{figures/guess3_changed_citrine_58f05f6d.jpg}{58f05f6d, saturated: rose quartz 54.7\%, aventurine 18.0\%}
           {figures/guess4_changed_citrine_58f05f6d.jpg}{saturated and cut: citrine 99.6\%, amethyst 0.3\%}
\sidebyside{figures/guess3_changed_citrine_2788e4fa.jpg}{2788e4fa, saturated: rose quartz 33.0\%, amethyst 25.1\%}
           {figures/guess4_changed_citrine_2788e4fa.jpg}{saturated and cut: citrine 97.9\%, amethyst 0.9\%}

What the results say about the guess, in my words: These results are the first
in which my guess was completely right. I can therefore assume that cutting the
background and making the crystal the main object of the image is an important
element that the model considers in order to answer correctly. The answers on
the cut images are so accurate that it makes me wonder whether it plays a bigger
role than colour (considering it was done on saturated images).

\subsection*{Guess 5: the cut on its own, on all my photos}

My guess, in my words as I gave it: ``At this point, I think a final check should
be done, with a 5th guess. Guess 4 showed that making the crystal bigger, and
cutting the background was important for the model to guess, finally, the 3
citrine correctly. However, in those images the color also was saturated from
the previous guess. So it could be interesting to test the cutting generally on
the images. In particular, applying the cut to all of my original images. Since
the other crystals were guessed right, observing how only the cut affects the
results can be useful to understand how much cutting the background and making
the crystal the main object in the image is taken in consideration from the
model. Also, because the images are original, so with normal saturation,
observing the results of the 3 citrine, can give an idea of whether cut plays a
bigger role than color. So my 5th guess is that if we cut the background from my
original images, making the crystals the main object in the image, the 3 citrine
6dc1dd3c-0216-45b8-abb7-214d633d280e.JPG, 58f05f6d-77f2-467e-b497-c4265bb43b8e.JPG,
and 2788e4fa-a7ad-453a-b3ad-3984b54d6e2c.JPG would still be guessed correctly,
and all the other images (that were previously guessed right) would also be
guessed right and more accurately.''

The change: all 21 of my original photos, with their normal colour, cut by the
same rule as Guess 4: a square centred on the stone, or on the group of stones
when there are several, its side 1.5 times the longer side of the stone, moved
inside the picture where it would stick out. The agent showed me all 21 cuts
before they were made, and on my request moved the square for
\texttt{amethyst/a95f9f93} lower, so that the stone is not cut.

My request to the agent, word for word:
\begin{quote}
So make a copy of each image in data/new\_images cut the background around the
crystals and save the copies in data/changed/cut
\end{quote}
and, after the preview:
\begin{quote}
yes, just change the cut on image a95f9f93-5117-4b70-9e92-4c513bb5decc.JPG to
make the cut a bit lower so that the crystal is not cut.
\end{quote}
The copies are in \texttt{data/changed/cut/<class>/}, one folder per class.
The agent wrote the cut as a cell that I ran on Colab. The page's first two
answers, the originals as the page answered them (amethyst and aventurine from
Guess 1, citrine from Guess 3, rose quartz read for this guess):

\begin{center}\scriptsize
\begin{tabular}{lll}
\toprule
photo & original & cut \\
\midrule
\texttt{amethyst/44db3aaf} & amethyst 96.7\%, rose quartz 2.4\% & amethyst 87.0\%, rose quartz 10.5\% \\
\texttt{amethyst/a95f9f93} & amethyst 86.9\%, rose quartz 8.1\% & amethyst 93.3\%, rose quartz 3.7\% \\
\texttt{amethyst/IMG\_3072} & amethyst 80.2\%, rose quartz 16.7\% & amethyst 62.1\%, rose quartz 29.2\% \\
\texttt{amethyst/IMG\_3089} & amethyst 90.5\%, aventurine 5.2\% & amethyst 88.7\%, aventurine 8.9\% \\
\texttt{amethyst/IMG\_3200} & amethyst 94.8\%, rose quartz 2.3\% & amethyst 91.1\%, rose quartz 5.7\% \\
\texttt{aventurine/59ad40f0} & aventurine 99.3\%, amethyst 0.4\% & aventurine 99.9\%, rose quartz 0.1\% \\
\texttt{aventurine/aba4d635} & aventurine 97.3\%, amethyst 1.7\% & aventurine 98.8\%, amethyst 0.7\% \\
\texttt{aventurine/IMG\_3068} & aventurine 75.2\%, rose quartz 15.6\% & aventurine 95.9\%, rose quartz 3.3\% \\
\texttt{aventurine/IMG\_3196} & aventurine 80.9\%, rose quartz 13.3\% & aventurine 98.8\%, rose quartz 0.7\% \\
\texttt{aventurine/IMG\_3197} & aventurine 88.7\%, rose quartz 9.9\% & aventurine 98.5\%, rose quartz 1.4\% \\
\texttt{citrine/6dc1dd3c} & rose quartz 52.8\%, aventurine 26.5\% & citrine 48.0\%, rose quartz 23.5\% \\
\texttt{citrine/08f618d6} & citrine 89.5\%, aventurine 4.7\% & citrine 94.5\%, aventurine 2.4\% \\
\texttt{citrine/58f05f6d} & rose quartz 49.1\%, aventurine 22.4\% & citrine 97.2\%, amethyst 2.2\% \\
\texttt{citrine/2788e4fa} & rose quartz 31.0\%, amethyst 29.8\% & citrine 90.1\%, amethyst 6.4\% \\
\texttt{citrine/f16efa23} & citrine 59.9\%, amethyst 25.1\% & citrine 85.9\%, aventurine 6.8\% \\
\texttt{rose\_quartz/9dacd40a} & rose quartz 99.8\%, aventurine 0.1\% & rose quartz 99.9\%, aventurine 0.1\% \\
\texttt{rose\_quartz/2270e0ff} & rose quartz 95.3\%, aventurine 2.9\% & rose quartz 97.6\%, aventurine 1.8\% \\
\texttt{rose\_quartz/ea7f4b76} & rose quartz 99.7\%, amethyst 0.2\% & rose quartz 99.8\%, aventurine 0.1\% \\
\texttt{rose\_quartz/IMG\_3069} & rose quartz 47.0\%, amethyst 34.8\% & rose quartz 53.3\%, aventurine 31.9\% \\
\texttt{rose\_quartz/IMG\_3194} & rose quartz 91.4\%, aventurine 6.9\% & rose quartz 99.1\%, aventurine 0.8\% \\
\texttt{rose\_quartz/IMG\_3204} & rose quartz 76.7\%, aventurine 22.2\% & rose quartz 63.2\%, aventurine 35.7\% \\
\bottomrule
\end{tabular}
\end{center}

The originals are 18 of 21 right on the page; the cut pictures are 21 of 21
right. The three citrines that were wrong are right without any extra colour,
\texttt{6dc1dd3c} at 48.0\% against 75.3\% when it was also saturated in Guess 4.
Of the 18 photos that were right before, 13 are more sure after the cut and 5
are less sure: four amethysts (\texttt{44db3aaf}, \texttt{IMG\_3072},
\texttt{IMG\_3089}, \texttt{IMG\_3200}) and one rose quartz (\texttt{IMG\_3204}).

\sidebyside{figures/guess5_original_amethyst_44db3aaf.jpg}{amethyst/44db3aaf, original: amethyst 96.7\%, rose quartz 2.4\%}
           {figures/guess5_changed_amethyst_44db3aaf.jpg}{cut: amethyst 87.0\%, rose quartz 10.5\%}
\sidebyside{figures/guess5_original_amethyst_a95f9f93.jpg}{amethyst/a95f9f93, original: amethyst 86.9\%, rose quartz 8.1\%}
           {figures/guess5_changed_amethyst_a95f9f93.jpg}{cut: amethyst 93.3\%, rose quartz 3.7\%}
\sidebyside{figures/guess5_original_amethyst_IMG_3072.jpg}{amethyst/IMG\_3072, original: amethyst 80.2\%, rose quartz 16.7\%}
           {figures/guess5_changed_amethyst_IMG_3072.jpg}{cut: amethyst 62.1\%, rose quartz 29.2\%}
\sidebyside{figures/guess5_original_amethyst_IMG_3089.jpg}{amethyst/IMG\_3089, original: amethyst 90.5\%, aventurine 5.2\%}
           {figures/guess5_changed_amethyst_IMG_3089.jpg}{cut: amethyst 88.7\%, aventurine 8.9\%}
\sidebyside{figures/guess5_original_amethyst_IMG_3200.jpg}{amethyst/IMG\_3200, original: amethyst 94.8\%, rose quartz 2.3\%}
           {figures/guess5_changed_amethyst_IMG_3200.jpg}{cut: amethyst 91.1\%, rose quartz 5.7\%}
\sidebyside{figures/guess5_original_aventurine_59ad40f0.jpg}{aventurine/59ad40f0, original: aventurine 99.3\%, amethyst 0.4\%}
           {figures/guess5_changed_aventurine_59ad40f0.jpg}{cut: aventurine 99.9\%, rose quartz 0.1\%}
\sidebyside{figures/guess5_original_aventurine_aba4d635.jpg}{aventurine/aba4d635, original: aventurine 97.3\%, amethyst 1.7\%}
           {figures/guess5_changed_aventurine_aba4d635.jpg}{cut: aventurine 98.8\%, amethyst 0.7\%}
\sidebyside{figures/guess5_original_aventurine_IMG_3068.jpg}{aventurine/IMG\_3068, original: aventurine 75.2\%, rose quartz 15.6\%}
           {figures/guess5_changed_aventurine_IMG_3068.jpg}{cut: aventurine 95.9\%, rose quartz 3.3\%}
\sidebyside{figures/guess5_original_aventurine_IMG_3196.jpg}{aventurine/IMG\_3196, original: aventurine 80.9\%, rose quartz 13.3\%}
           {figures/guess5_changed_aventurine_IMG_3196.jpg}{cut: aventurine 98.8\%, rose quartz 0.7\%}
\sidebyside{figures/guess5_original_aventurine_IMG_3197.jpg}{aventurine/IMG\_3197, original: aventurine 88.7\%, rose quartz 9.9\%}
           {figures/guess5_changed_aventurine_IMG_3197.jpg}{cut: aventurine 98.5\%, rose quartz 1.4\%}
\sidebyside{figures/guess5_original_citrine_6dc1dd3c.jpg}{citrine/6dc1dd3c, original: rose quartz 52.8\%, aventurine 26.5\%}
           {figures/guess5_changed_citrine_6dc1dd3c.jpg}{cut: citrine 48.0\%, rose quartz 23.5\%}
\sidebyside{figures/guess5_original_citrine_08f618d6.jpg}{citrine/08f618d6, original: citrine 89.5\%, aventurine 4.7\%}
           {figures/guess5_changed_citrine_08f618d6.jpg}{cut: citrine 94.5\%, aventurine 2.4\%}
\sidebyside{figures/guess5_original_citrine_58f05f6d.jpg}{citrine/58f05f6d, original: rose quartz 49.1\%, aventurine 22.4\%}
           {figures/guess5_changed_citrine_58f05f6d.jpg}{cut: citrine 97.2\%, amethyst 2.2\%}
\sidebyside{figures/guess5_original_citrine_2788e4fa.jpg}{citrine/2788e4fa, original: rose quartz 31.0\%, amethyst 29.8\%}
           {figures/guess5_changed_citrine_2788e4fa.jpg}{cut: citrine 90.1\%, amethyst 6.4\%}
\sidebyside{figures/guess5_original_citrine_f16efa23.jpg}{citrine/f16efa23, original: citrine 59.9\%, amethyst 25.1\%}
           {figures/guess5_changed_citrine_f16efa23.jpg}{cut: citrine 85.9\%, aventurine 6.8\%}
\sidebyside{figures/guess5_original_rose_quartz_9dacd40a.jpg}{rose quartz/9dacd40a, original: rose quartz 99.8\%, aventurine 0.1\%}
           {figures/guess5_changed_rose_quartz_9dacd40a.jpg}{cut: rose quartz 99.9\%, aventurine 0.1\%}
\sidebyside{figures/guess5_original_rose_quartz_2270e0ff.jpg}{rose quartz/2270e0ff, original: rose quartz 95.3\%, aventurine 2.9\%}
           {figures/guess5_changed_rose_quartz_2270e0ff.jpg}{cut: rose quartz 97.6\%, aventurine 1.8\%}
\sidebyside{figures/guess5_original_rose_quartz_ea7f4b76.jpg}{rose quartz/ea7f4b76, original: rose quartz 99.7\%, amethyst 0.2\%}
           {figures/guess5_changed_rose_quartz_ea7f4b76.jpg}{cut: rose quartz 99.8\%, aventurine 0.1\%}
\sidebyside{figures/guess5_original_rose_quartz_IMG_3069.jpg}{rose quartz/IMG\_3069, original: rose quartz 47.0\%, amethyst 34.8\%}
           {figures/guess5_changed_rose_quartz_IMG_3069.jpg}{cut: rose quartz 53.3\%, aventurine 31.9\%}
\sidebyside{figures/guess5_original_rose_quartz_IMG_3194.jpg}{rose quartz/IMG\_3194, original: rose quartz 91.4\%, aventurine 6.9\%}
           {figures/guess5_changed_rose_quartz_IMG_3194.jpg}{cut: rose quartz 99.1\%, aventurine 0.8\%}
\sidebyside{figures/guess5_original_rose_quartz_IMG_3204.jpg}{rose quartz/IMG\_3204, original: rose quartz 76.7\%, aventurine 22.2\%}
           {figures/guess5_changed_rose_quartz_IMG_3204.jpg}{cut: rose quartz 63.2\%, aventurine 35.7\%}

What the results say about the guess, in my words: From these results I can
observe two things. First, as I expected, the 3 wrong citrines were guessed
correctly even without being more saturated, although with less accuracy. This
means that cutting the background and making the crystal bigger in the image is
important for the model to guess it correctly; in particular it plays a bigger
role than colour, since in Guess 3 the 3 citrines were still guessed wrong when
only the saturation was enhanced. Second, the crystals that were already guessed
right were also guessed right, as expected, but not exactly as I imagined: not
all the answers were more accurate. In particular, 4 amethysts were guessed with
less accuracy and only one with more. This is important to give the right amount
of importance to the cut (so to the size of the crystals and the absence of
background). Considering the other factors that I observed in the previous
guesses, such as transparency and colour, and my idea of a combination of
elements, the reason why the 4 amethysts got a less accurate answer after the
cut could be connected to the fact that the colour of amethyst is fuller and
darker, which can have an effect on how useful the cut is to the model in that
case. However, it is important to also mention that one of the rose quartz
photos, IMG\_3204.jpg, also had less accuracy after the cut, but in this case
the transparency and saturation of the crystal are different from amethyst.
Therefore, it is possible that there is another element that the model looks at
that is not among those I tried to test in my 5 guesses.

\subsection*{What I now think the model uses}
At the beginning of the assignment I thought that the model would mostly look at
the colour of the crystals to distinguish them, as it is the first thing that is
visible by looking at them. However, through the 5 guesses I made I became more
sure that the model does not use only one element to recognise the different
crystals, but rather takes a combination of them into consideration. This is
probably the part of the conclusion that I am most sure about.

The elements that I tried to test, in order to understand whether they are
considered by the model, were three. The first was colour, so the saturation of
the crystals. The second was transparency, in the sense of how pale the crystal
is, and it was observed through less contrast. Lastly, the size of the crystal
and the reduction of the background, which were observed together through the
cut. Among these, I think that mainly two elements are considered by the model,
but differently. The element that I think has more importance for the model is
the size of the crystal and the reduction of the background, since all my images
were recognised correctly after the cut that made the crystal the main object in
the image and reduced the background. Colour (saturation) is the second element
in importance, considering that the 3 saturated citrines were not recognised
correctly before the cut. As for transparency, it might play a role, but it was
not easy to check, as lowering the contrast is not exactly the same as having a
crystal that is pale and less full in colour. Therefore, from my guesses I
conclude that the model uses a combination of the size of the crystal, the
reduction of the background, and colour to distinguish the categories, and that
the size of the crystal and the reduction of the background have a stronger
impact on the model than colour.

However, it is important to mention that my guesses were based on tests I tried
on my own images, which were only 21, so my conclusion has limitations. As
mentioned in the conclusion of Guess 5, I think there might be other elements
that the model looks at to distinguish the categories, which I did not test in my
guesses. Apart from Guess 5, which tested the cut on all my images, my other
guesses were done on a smaller number of images; in particular, my second and
fourth guesses were more specific and focused only on 3 amethysts and on the 3
wrong citrines, so my conclusions from those two guesses are limited. I am aware
that more tests and guesses, especially on more images and from other sources as
well, could give a more complete answer to what the model uses to recognise the
categories.

% ---------------------------------------------------------------------------
\section*{Appendix. Output of check.py}
\outputfile{outputs/check.txt}

\end{document}
